% Avoid simply summarizing the findings of each paper, one paragraph per paper. \cite{hosseini2016adaptive} discusses.. \cite{SURV} writes  

% Instead, synthesize the papers by identifying themes and using the papers as references for the themes. 

% NEED TO have conference co chairs if citing a conference paper

% journals from the last 2-3 years
The architectural approach of this research is heavily informed by previous exploration into hardware-software co-design for implantable medical devices. Historically, the fundamental challenge of respecting the strict $1^\circ$C thermal limit has forced designers to rely on highly specialized, monolithic ASICs. Kassiri et al.'s closed-loop neurostimulator \cite{7982670}, Neuralink's N1 ASIC \cite{Musk2019}, and the NeuroPace RNS System \cite{Heck2014} represent the state-of-the-art implementations of this rigid baseline. These systems achieve unparalleled energy efficiency for their exact medical or recording tasks by permanently baking clinical logic into the physical silicon. However, this comes at the cost of algorithmic flexibility. Unlike these ASICs, HALO is built on an open-source instruction set architecture (ISA), RISC-V Vector extension (RVV). This enables its binaries to pick up improvements from every generation of the open-source GCC backend without a tape-out. The Results chapter demonstrates the payoff: simply updating the compiler cut collective execution cost by 7.2~percent across the core arithmetic and logical operations (excluding Octave-mode comparison tests), and in three cases (\texttt{and} dbl$\times$dbl, \texttt{and} cpx$\times$dbl, \texttt{or} cpx$\times$dbl) the auto-vectorized scalar code now beats the hand-written intrinsics outright.

% What cannot be changed post-tape-out is the accelerator datapath, the register file width, and the instruction set available to that microcontroller. It is that second layer---the instruction set and the datapath that executes it---where HALO's RISC-V plus RVV choice matters. Because the RVV encoding is fixed but the micro-architecture behind it is not, and because RVV is vector-length agnostic, the same compiled binary picks up improvements from every generation of the open-source GCC backend without a tape-out. The Results chapter demonstrates the payoff: simply updating the compiler cut collective execution cost by 7.2~percent across the core arithmetic and logical operations (excluding Octave-mode comparison tests), and in three cases (\texttt{and} dbl$\times$dbl, \texttt{and} cpx$\times$dbl, \texttt{or} cpx$\times$dbl) the auto-vectorized scalar code now beats the hand-written intrinsics outright.

\section{Translation for High-Level Languages}

In the domain of high-level synthesis and automated code generation, commercial tools like MathWorks' MATLAB Coder \cite{mathworks_matlab_coder} provide a highly refined reference point. MATLAB Coder is highly capable at parsing proprietary MATLAB scripts and generating optimized C and C++ code targeted primarily at advanced desktop processors and commercial ARM architectures. Similarly, for researchers utilizing Python—another dynamically typed language widely used in BCI data analysis and real-time processing pipelines—tools like Cython \cite{behnel2011cython} offer similar benefits. Cython addresses the performance bottlenecks of vanilla Python by compiling code directly into efficient C or C++ extensions, utilizing static type definitions to achieve substantial execution speedups for data pipelines.

Although both MATLAB Coder and Cython are highly effective for accelerating signal processing on standard architectures, they share a critical limitation in the context of this research: they are structurally agnostic regarding emerging, open-source instruction sets. Neither tool inherently supports the automated generation of targeted RISC-V Vector (RVV) intrinsics required to meet the ultra-low-power constraints of the HALO platform. Consequently, the custom ts-traversal and RISCV-Matrix integration developed in this work is necessary to bridge the semantic gap between high-level BCI algorithm prototyping and the specific vector execution demands of the embedded hardware.
